\section*{Theory}

The SIR Model is derived from a population of fixed size $N$ that is partitioned at every time $t$ into categories:
\newcommand{\sym}[1]{\makebox[1.1em][r]{$#1$}}
\begin{itemize}
  \item \sym{S} : \textbf{S}usceptible individuals,
  \item \sym{I} : \textbf{I}nfectious individuals,
  \item \sym{R} : \textbf{R}ecovered (or removed/dead) individuals.
\end{itemize}
Here $S=S(t)$, $I=I(t)$, $R=R(t)$ denote the sizes of the categories. Because every individual occupies exactly one category and the total population is
conserved, the state satisfies
\begin{equation}
  I_0 + S_0 \;=\; I + S + R \;=\; N ,
  \label{eq:conservation}
\end{equation}
where $N$ is the total population.
\newline
\paragraph{Initial conditions.}
The model is seeded with
\begin{equation}
  I_0 = \#\{\text{infected at } t=0\}, \qquad
  S_0 = N - I_0, \qquad
  R_0^{\,(\text{init})} = 0 ,
  \label{eq:ic}
\end{equation}
so that the conservation law \eqref{eq:conservation} holds from $t=0$
onward.\footnote{We write $R_0^{\,(\text{init})}$ for the initial number of
recovered individuals (which is $0$) to avoid clashing with the basic
reproduction number $R_0=\beta/\gamma$ introduced in
Section~\ref{sec:ode}.}
\newline
\paragraph{Transition rates}
Two mechanisms drive the dynamics of the system. The \emph{infection} (arrival into $I$) and
\emph{recovery} (departure from $I$) rates are
\begin{equation}
  \lambda_{\mathrm{inf}}(S,I) = \frac{\beta}{N}\,S\,I,
  \qquad
  \lambda_{\mathrm{rec}}(I) = \gamma\, I,
  \qquad \beta,\gamma \in \mathbb{R}^{+}.
  \label{eq:rates}
\end{equation}
The parameters are interpreted as follows. $\beta$ measures
\emph{how infectious an individual is} (the effective contact/transmission
rate); $\gamma$ is the \emph{recovery rate}. Indeed, if r.v. $A$ denotes the time an individual remains infectious,
\begin{equation}
  A \sim \mathrm{Exp}(\gamma)
  \quad\Longrightarrow\quad
  \mathbb{E}[A] = \frac{1}{\gamma},
  \label{eq:meaninf}
\end{equation}
so $1/\gamma$ is the \emph{average infectious time}.

\section{Deterministic dynamics, SIR modelled from system of ODEs}
\label{sec:ode}

Replacing the random transitions by their rates gives the deterministic SIR
system of ordinary differential equations
\begin{equation}
  \left\{
  \begin{aligned}
    \frac{dS}{dt} &= -\frac{\beta\,S\,I}{N}
                   &&= -\,\lambda_{\mathrm{inf}}(S,I),\\[4pt]
    \frac{dI}{dt} &= \frac{\beta\,S\,I}{N} - \gamma I
                   &&= \;\lambda_{\mathrm{inf}}(S,I) - \lambda_{\mathrm{rec}}(I),\\[4pt]
    \frac{dR}{dt} &= \gamma I
                   &&= \;\lambda_{\mathrm{rec}}(I).
  \end{aligned}
  \right.
  \label{eq:ode}
\end{equation}
The three equations sum to $\tfrac{d}{dt}(S+I+R)=0$, consistent with the
conservation law \eqref{eq:conservation}.
\newline
\paragraph{Basic reproduction number}
The dimensionless ratio
\begin{equation}
  R_0 = \frac{\beta}{\gamma}
  \label{eq:R0}
\end{equation}
controls the behaviour of the model: \emph{if $R_0 > 1$ an epidemic occurs}
(the infectious compartment can grow initially), whereas $R_0<1$ leads to
immediate die-out.

\section{Going stochastic: SIR as a stochastic process}

For finite $N$ the category counts are integer valued and the dynamics are
stochastic. We model the transitions as a \emph{Poisson arrival process}:
individuals \emph{arrive} into $I$ (infection) and \emph{depart} from $I$
(recovery). The nice thing about this is that if you now extend the graph we can now simply update the stochastic system \ref{eq:stoch_balance}.

\begin{center}
\begin{tikzpicture}[
    node distance=22mm and 26mm,
    comp/.style={circle, draw, minimum size=11mm, font=\large},
    every edge/.style={draw, -{Latex[length=2.5mm]}}
  ]
  \node[comp] (S) {$S$};
  \node[comp, right=of S] (I) {$I$};
  \node[comp, right=of I] (R) {$R$};

  \draw (S) edge node[above]{Arrival} node[below]{$\lambda_{\mathrm{inf}}$} (I);
  \draw (I) edge node[above]{Departure} node[below]{$\lambda_{\mathrm{rec}}$} (R);
\end{tikzpicture}
\end{center}

The same two rates \eqref{eq:rates} govern the process, and the (now
operator-level) balance equations read
\begin{equation}
  \left\{
  \begin{aligned}
    \frac{dS}{dt} &= -\,\lambda_{\mathrm{inf}}(I,S),\\[2pt]
    \frac{dI}{dt} &= \;\lambda_{\mathrm{inf}}(S,I) - \lambda_{\mathrm{rec}}(I),\\[2pt]
    \frac{dR}{dt} &= \;\lambda_{\mathrm{rec}}(I).
  \end{aligned}
  \right.
  \label{eq:stoch_balance}
\end{equation}
The process is a continuous-time Markov chain on the state space
$\{(S,I,R)\in\mathbb{Z}_{\ge 0}^3 : S+I+R=N\}$, with only two admissible
jumps from a state $(S,I,R)$:
\begin{align}
  \text{infection:} &\quad (S,I,R)\;\longrightarrow\;(S-1,\,I+1,\,R),
    && \text{rate } \lambda_{\mathrm{inf}}(S,I),\\
  \text{recovery:}  &\quad (S,I,R)\;\longrightarrow\;(S,\,I-1,\,R+1),
    && \text{rate } \lambda_{\mathrm{rec}}(I).
\end{align}

\section{Embedding in the Markov chain: competing exponential clocks}
\label{sec:embed}
\footnote{These derivations are based on the slides regarding markov chains, \eqref{eq:PEA_int} was expanded for understanding.}

Associate to the two possible events independent exponential random variables: an
\emph{arrival} (infection) r.v. $A$ and a \emph{departure} (recovery) r.v.
$D$,
\begin{equation}
  A \sim \mathrm{Exp}\!\big(\lambda_{\mathrm{inf}}(S,I)\big),
  \qquad
  D \sim \mathrm{Exp}\!\big(\lambda_{\mathrm{rec}}(I)\big).
\end{equation}
The next event occurs at the first r.v. to be observed,
\begin{equation}
  E = \min(A,D),
\end{equation}
and the embedding answers two questions: \emph{when} does the next event
happen (the holding time $E$), and \emph{which} event is it ($E=A$ vs.\ $E=D$).

\subsection{Holding time distribution}

Using independence of $A$ and $D$,
\begin{equation}
  P(E>t) = P(A>t,\,D>t) = P(A>t)\,P(D>t)
         = e^{-\lambda_{\mathrm{inf}}t}\;e^{-\lambda_{\mathrm{rec}}t}
         = e^{-(\lambda_{\mathrm{inf}}+\lambda_{\mathrm{rec}})t}.
\end{equation}
Hence the minimum of the two exponential clocks is again exponential, with rate
equal to the \emph{sum} of the individual rates:
\begin{equation}
  E \sim \mathrm{Exp}\!\big(\lambda_{\mathrm{inf}}+\lambda_{\mathrm{rec}}\big)
    = \mathrm{Exp}\!\left(\Big(\tfrac{\beta}{N}S+\gamma\Big)I\right),
  \label{eq:holding}
\end{equation}
since
$\lambda_{\mathrm{inf}}+\lambda_{\mathrm{rec}}
 = \tfrac{\beta}{N}SI+\gamma I = \big(\tfrac{\beta}{N}S+\gamma\big)I$.

\subsection{Which event occurs}

The probability that the next event is an infection is $P(E=A)=P(A<D)$.
Conditioning on the arrival time and summing over the infinitesimal intervals
$[t,t+dt]$,
\begin{align}
  P(E=A) = P(A<D)
  &= \sum P\big(A<D,\; A\in[t,t+dt]\big) \notag\\
  &= \sum P\big(A<D \mid A\in[t,t+dt]\big)\,P\big(A\in[t,t+dt]\big) \notag\\
  &\xrightarrow[dt\to 0]{}
     \int_{0}^{\infty} P\big(A<D \mid A=t\big)\,f_A(t)\,dt .
  \label{eq:PEA_int}
\end{align}
By independence, $P(A<D\mid A=t)=P(D>t)=e^{-\lambda_{\mathrm{rec}}t}$, and
$f_A(t)=\lambda_{\mathrm{inf}}e^{-\lambda_{\mathrm{inf}}t}$. The integral is the
standard competing-exponentials computation. Writing it with the board's
generic rate symbols $\lambda$ (arrival) and $k\mu$ (departure),
\begin{equation}
  \int_{0}^{\infty} e^{-k\mu t}\,\lambda e^{-\lambda t}\,dt
  = \int_{0}^{\infty} \lambda\, e^{-(\lambda+k\mu)t}\,dt
  = \frac{\lambda}{\lambda + k\mu}.
  \label{eq:PEA_generic}
\end{equation}
Substituting the SIR rates
\begin{equation}
  \lambda = \lambda_{\mathrm{inf}} = \frac{\beta SI}{N},
  \qquad
  k\mu = \lambda_{\mathrm{rec}} = \gamma I,
\end{equation}
gives the explicit selection probability
\begin{equation}
  \boxed{\;
  P(E=A)
  = \frac{\beta SI/N}{\,\beta SI/N + \gamma I\,}
  = \frac{\lambda_{\mathrm{inf}}(I,S)}
         {\lambda_{\mathrm{inf}}(I,S)+\lambda_{\mathrm{rec}}(I)} \; }.
  \label{eq:PEA}
\end{equation}
Consequently $P(E=D)=1-P(E=A)=\dfrac{\gamma I}{\beta SI/N+\gamma I}$.

\section{The Gillespie algorithm}

Equations \eqref{eq:holding} and \eqref{eq:PEA} are now sufficient to do a stochastic
simulation. From a state $(S,I,R)$:

\begin{enumerate}
  \item \textbf{Define the random variables}
        \[
          A \sim \mathrm{Exp}\!\big(\lambda_{\mathrm{inf}}(S,I)\big),\qquad
          D \sim \mathrm{Exp}\!\big(\lambda_{\mathrm{rec}}(I)\big),\qquad
          E = \min(A,D).
        \]
  \item \textbf{Draw} $U \sim \mathrm{Uniform}(0,1)$.
  \item \textbf{Select the event} using \eqref{eq:PEA}:
        \[
          \begin{cases}
            \text{if } U \le P(E=A): &
              (S,I,R)\;\longrightarrow\;(S-1,\,I+1,\,R)
              \quad\text{(infection)},\\[4pt]
            \text{else } \big(U > P(E=A)\big): &
              (S,I,R)\;\longrightarrow\;(S,\,I-1,\,R+1)
              \quad\text{(recovery)}.
          \end{cases}
        \]
  \item \textbf{Advance time} by the holding time, $t \leftarrow t + E$, and
        repeat until $I=0$ (no further events are possible).
\end{enumerate}

Notice if new rates are introduce to model differently it would instead of being a binary decision be chosen for categorical distribution. 